% !TeX root = main.tex
\changecolours{ScoutTeal}
\section[Agrupamento Analítico e Acumuladores]{Agrupamento Analítico e Operadores Acumuladores}

% ------------------------------------------------------------------------------
\begin{frame}[fragile]{O Estágio \texttt{\$group}: A Chave de Agrupamento \texttt{\_id}}
	\vspace{0.1cm}
	\begin{columns}[t,onlytextwidth]
		
		% Coluna 1
		\begin{column}{0.48\textwidth}
			\alert{Mecanismo de Agrupamento}\par
			\footnotesize
			\begin{itemize}
				\item Reduz múltiplos documentos em um único documento por valor distinto de \lstinline[language=mongosh, basicstyle=\footnotesize\ttfamily]|_id|.
				
				\item \textbf{Agrupamento Global (KPI):} Usar \lstinline[language=mongosh, basicstyle=\footnotesize\ttfamily]|_id: null| agrega toda a coleção em um único registro.
				
				\item \textbf{Chaves Compostas:} Agrupa por múltiplos atributos via subdocumento:
				
				\vspace{0.15cm}
				\centering
				\lstinline[language=mongosh, basicstyle=\scriptsize\ttfamily]|_id: { ano: "$ano", mes: "$mes" }|
			\end{itemize}
		\end{column}
		
		% Coluna 2
		\begin{column}{0.48\textwidth}
			\alert{Atenção: A Regra do Cifrão (\texttt{\$})}\par
			\footnotesize
			\begin{itemize}
				\item Para acessar o \textbf{valor de um campo}, é obrigatório usar o prefixo \texttt{\$}:
				
				\vspace{0.15cm}
				\centering
				\lstinline[language=mongosh, basicstyle=\footnotesize\ttfamily]|{ $sum: "$pagamento.total_pago" }|
				
				\vspace{0.15cm}
				\raggedright
				\item Sem o prefixo \texttt{\$}, interpreta-se o nome como texto literal constante:
				
				\vspace{0.15cm}
				\centering
				\lstinline[language=mongosh, basicstyle=\footnotesize\ttfamily]|{ $sum: "pagamento.total_pago" }|
				\par\vspace{0.1cm}
				$\implies$ \textcolor{red!80!black}{\textbf{Erro / Soma Zero!}}
			\end{itemize}
		\end{column}
		
	\end{columns}
\end{frame}

% ------------------------------------------------------------------------------
\begin{frame}[fragile]{Quadro de Acumuladores Fundamentais}
  \vspace{0.05cm}
  \begin{columns}[t,onlytextwidth]
    \begin{column}{0.45\textwidth}
      \alert{Acumuladores Mais Utilizados}\par
      \vspace{0.4cm}
      \scriptsize
      \begin{tabular}{@{}lp{3.8cm}@{}}
        \toprule
        \textbf{Operador} & \textbf{Finalidade Analítica} \\
        \midrule
        \lstinline[language=mongosh, basicstyle=\footnotesize\ttfamily]|$sum| & Soma ou contagem de itens \\
        \lstinline[language=mongosh, basicstyle=\footnotesize\ttfamily]|$avg| & Média aritmética dos campos \\
        \lstinline[language=mongosh, basicstyle=\footnotesize\ttfamily]|$min/$max| & Menor e maior valor \\
        \lstinline[language=mongosh, basicstyle=\footnotesize\ttfamily]|$first/$last| & Primeiro e último por ordem \\
        \lstinline[language=mongosh, basicstyle=\footnotesize\ttfamily]|$push| & Array com histórico \\
        \lstinline[language=mongosh, basicstyle=\footnotesize\ttfamily]|$addToSet| & Conjunto sem duplicatas \\
        \bottomrule
      \end{tabular}
    \end{column}
    \begin{column}{0.51\textwidth}
      \alert{Diferença: \texttt{\$push} vs \texttt{\$addToSet}}\par
    
      \begin{itemize}
        \item \textbf{\texttt{\$push}:} Mantém todas as ocorrências de itens, preservando repetições. Indicado para rastreabilidade e histórico de pedidos.
        \item \textbf{\texttt{\$addToSet}:} Opera como conjunto matemático estrito, eliminando valores duplicados. Indicado para listar categorias únicas compradas.
      \end{itemize}
    \end{column}
  \end{columns}
\end{frame}

% ------------------------------------------------------------------------------
\begin{frame}[fragile]{Exemplo Pareado: Análise Mensal de Faturamento}
  \vspace{0.05cm}
  \begin{columns}[t,onlytextwidth]
    \begin{column}{0.44\textwidth}
    	\alert{Formulação Matemática}\par
    	\vspace{0.2cm}
    	\scriptsize
    	Seja o lote de pedidos entregues em Outubro:
    	\begin{center}
    		$P = \{150.00, 320.00, 480.00, 50.00\}$
    	\end{center}
    	\begin{itemize}
    		\item \textbf{Faturamento Total ($\text{Fat}$):}
    		\begin{itemize}
    			\item[] $\text{Fat} = \sum_{i=1}^{4} p_i = 1000.00$
    		\end{itemize}
    		\item \textbf{Ticket Médio ($\bar{X}$):}
    		\begin{itemize}
    			\item[] $\bar{X} = \frac{1000.00}{4} = 250.00$
    		\end{itemize}
    		\item \textbf{Volume Total de Pedidos:}
    		\begin{itemize}
    			\item[] $N = \sum 1 = 4$
    		\end{itemize}
    	\end{itemize}
    \end{column}
    \begin{column}{0.54\textwidth}
      \alert{Implementação Equivalente no MongoSH}\par
      \vspace{0.05cm}
      \begin{lstlisting}[language=mongosh, basicstyle=\scriptsize\ttfamily]
db.pedidos.aggregate([
  { $match: { status: "ENTREGUE" } },
  {
    $group: {
      _id: {
        ano: { $year: "$data_criacao" },
        mes: { $month: "$data_criacao" }
      },
      faturamento: { $sum: "$pagamento.total_pago" },
      ticket_medio: { $avg: "$pagamento.total_pago" },
      total_pedidos: { $sum: 1 }
    }
  },
  { $sort: { "_id.ano": 1, "_id.mes": 1 } }
])
      \end{lstlisting}
    \end{column}
  \end{columns}
\end{frame}
